% Appendix to the manuscript in ../manuscript/main.tex: the frozen system
% prompt, the event texts, one complete model call, and the per-run counts
% for the 40 model runs. RA-L counts an appendix toward its page limit, so
% the paper cites this file by its repository URL. The text files beside
% this one are the prompt and the calls as the delivery crate wrote them.
% Build: tectonic -o <folder> main.tex
\documentclass[conference]{IEEEtran}
\usepackage[T1]{fontenc}
\usepackage{times}
\usepackage[varqu]{zi4}
\usepackage{booktabs}
\usepackage{tabularx}
\usepackage{listings}
\lstset{basicstyle=\footnotesize\ttfamily,breaklines=true,
  breakatwhitespace=true,
  columns=fullflexible,keepspaces=true,frame=single,framerule=0.2pt,
  xleftmargin=2pt,xrightmargin=2pt,aboveskip=4pt,belowskip=6pt}
\usepackage[bookmarks=true,hidelinks]{hyperref}

\title{Appendix to Ripple: Measuring the Cost of Over-Helping\\in Robot
Fleets}
\author{\IEEEauthorblockN{Ethan Moosbrugger Boos}
\IEEEauthorblockA{Georgia Institute of Technology}}

\begin{document}
\maketitle

This appendix holds the material the paper refers to but has no room
for: the system prompt every model call used, the one text for each event
kind, one complete live call and one veto prompt with their replies, and
the per-run counts for the 40 model runs. The study code is the delivery
crate of the Ripple repository, and the prompt and call texts here are the
files beside this document, as the crate wrote them.

\section{The frozen system prompt}

Every model call in the study used this system prompt. It was written once,
before any scored run, and never changed. The mission file holds it.

\lstinputlisting{system-prompt.txt}

\section{Event texts}
\label{app:events}

Each event kind has one text, written once and filled in from the event's
fields. The model reads events as these sentences under \emph{New events}
and \emph{Events still in force} (Section~\ref{app:call}). The placeholders
are the hub or hubs, the site, the stock item, the area, the vehicle class,
and the start and end minutes.

\begin{center}
\footnotesize
\begin{tabularx}{\columnwidth}{@{}lX@{}}
\toprule
Kind & Text \\
\midrule
Surge & A mass-casualty incident near \{hubs\}: the hospitals there are
sending trauma requests beyond what those hubs can fly alone, starting at
minute \{from\}. \\
Shortage & \{item\} has run out at \{hubs\} as of minute \{from\}. Requests
for \{item\} at their sites can only be served from hubs that still hold
it. \\
Closure & Airspace is closed in \{area\} from minute \{from\} until minute
\{to\}; no flights through it. \\
Outage & \{hubs\} cannot launch from minute \{from\} until minute \{to\};
its sites must be covered by other hubs. \\
Weather & Weather at \{hubs\}: no \{class\} launches from minute \{from\}
until minute \{to\}. \\
Priority & Hospital \{site\} takes priority over routine deliveries from
minute \{from\} until \{to\}. \\
Rooftop closed & The rooftop at hospital \{site\} is closed from minute
\{from\} until minute \{to\}; only drops at its ground pad are allowed. \\
Corridor & A medevac helicopter corridor is active in \{area\} from minute
\{from\} until minute \{to\}; no vehicle may cross it below 150 m. \\
\bottomrule
\end{tabularx}
\end{center}

\section{One complete model call}
\label{app:call}

The prompt below is one call to Opus as a live manager with the exceptions
view and feedback on a 16-hub day, at minute 35, and the reply that came
back. The prompt is what the delivery crate wrote; nothing was removed. The
order format at the end is part of every live prompt. The reply closed the
rooftop that had just been reported and moved two fixed-wing aircraft from
hub 5 to hub 4, the one hub with waiting requests and no idle vehicle, with
a return time.

\lstinputlisting{call-prompt.txt}

The reply:

\lstinputlisting{call-reply.txt}

A veto prompt has the same head and, instead of the order format, the
replanner's numbered proposals and a one-line instruction. The tail of one
such prompt, at minute 41 on a 16-hub day after a stock shortage at five
hubs, and its reply:

\lstinputlisting{veto-tail.txt}

\lstinputlisting{veto-reply.txt}

\section{Per-run counts}
\label{app:counts}

Table~\ref{tab:counts} gives, for each of the 40 model runs in
the paper's model table, the number of model calls, the orders the model
gave (or, for the veto, the replanner's proposals it kept), how many were
applied and how many refused by the dispatcher, the replies that could not
be read, the moves (applied relocations and stock flights), the lending
permissions, and the vehicles lost.

\begin{table*}[h]
\caption{Per-run counts for the 40 model runs. Refused orders are those
the dispatcher could not apply, almost all for too few idle vehicles.
Unreadable replies were not JSON in the expected form and changed nothing.}
\label{tab:counts}
\centering
\footnotesize
\begin{tabular}{llllrrrrrrrr}
\toprule
Hubs & Role & Model & Seed & Calls & Orders & Applied & Refused & Unreadable & Moves & Lends & Lost \\
\midrule
16 & Live & Haiku & 1 & 49 & 94 & 90 & 4 & 0 & 7 & 45 & 5 \\
16 & Live & Haiku & 2 & 50 & 69 & 65 & 4 & 0 & 7 & 18 & 3 \\
16 & Live & Haiku & 3 & 53 & 120 & 117 & 3 & 1 & 10 & 70 & 0 \\
16 & Live & Sonnet & 1 & 49 & 171 & 144 & 27 & 2 & 36 & 76 & 4 \\
16 & Live & Sonnet & 2 & 50 & 137 & 129 & 8 & 1 & 48 & 47 & 9 \\
16 & Live & Sonnet & 3 & 53 & 183 & 169 & 14 & 0 & 56 & 80 & 12 \\
16 & Live & Opus & 1 & 49 & 385 & 326 & 59 & 0 & 85 & 192 & 5 \\
16 & Live & Opus & 2 & 50 & 382 & 336 & 46 & 0 & 71 & 189 & 10 \\
16 & Live & Opus & 3 & 53 & 409 & 379 & 30 & 0 & 54 & 250 & 4 \\
16 & Live & GPT-6 Luna & 1 & 49 & 242 & 235 & 7 & 0 & 2 & 193 & 0 \\
16 & Live & GPT-6 Luna & 2 & 50 & 235 & 226 & 9 & 1 & 6 & 164 & 1 \\
16 & Live & GPT-6 Luna & 3 & 53 & 266 & 260 & 6 & 1 & 1 & 183 & 0 \\
\addlinespace
16 & Veto & Haiku & 1 & 47 & 205 & 205 & 0 & 0 & 16 & 152 & 0 \\
16 & Veto & Haiku & 2 & 48 & 225 & 225 & 0 & 0 & 18 & 168 & 0 \\
16 & Veto & Haiku & 3 & 46 & 232 & 232 & 0 & 0 & 23 & 172 & 0 \\
16 & Veto & Sonnet & 1 & 46 & 244 & 244 & 0 & 0 & 19 & 184 & 0 \\
16 & Veto & Sonnet & 2 & 46 & 156 & 155 & 1 & 0 & 34 & 79 & 0 \\
16 & Veto & Sonnet & 3 & 46 & 152 & 151 & 1 & 0 & 23 & 88 & 0 \\
16 & Veto & Opus & 1 & 46 & 213 & 213 & 0 & 0 & 20 & 152 & 0 \\
16 & Veto & Opus & 2 & 46 & 145 & 144 & 1 & 0 & 30 & 72 & 0 \\
16 & Veto & Opus & 3 & 48 & 190 & 189 & 1 & 0 & 20 & 128 & 0 \\
16 & Veto & GPT-6 Luna & 1 & 47 & 289 & 289 & 0 & 0 & 18 & 232 & 0 \\
16 & Veto & GPT-6 Luna & 2 & 46 & 237 & 235 & 2 & 0 & 25 & 168 & 0 \\
16 & Veto & GPT-6 Luna & 3 & 46 & 315 & 314 & 1 & 0 & 23 & 250 & 0 \\
\addlinespace
64 & Live & Haiku & 1 & 185 & 362 & 335 & 27 & 0 & 61 & 121 & 34 \\
64 & Live & Haiku & 2 & 185 & 348 & 321 & 27 & 2 & 67 & 98 & 21 \\
64 & Live & Sonnet & 1 & 185 & 876 & 713 & 163 & 3 & 296 & 264 & 48 \\
64 & Live & Sonnet & 2 & 185 & 989 & 793 & 196 & 6 & 306 & 329 & 56 \\
64 & Live & Opus & 1 & 185 & 1843 & 1573 & 270 & 0 & 623 & 496 & 49 \\
64 & Live & Opus & 2 & 185 & 1892 & 1651 & 241 & 0 & 339 & 899 & 47 \\
64 & Live & GPT-6 Luna & 1 & 185 & 1577 & 1481 & 96 & 1 & 72 & 1096 & 11 \\
64 & Live & GPT-6 Luna & 2 & 185 & 1450 & 1394 & 56 & 2 & 39 & 1000 & 4 \\
\addlinespace
64 & Veto & Haiku & 1 & 184 & 1060 & 1041 & 19 & 0 & 134 & 774 & 0 \\
64 & Veto & Haiku & 2 & 184 & 1195 & 1165 & 30 & 0 & 169 & 848 & 0 \\
64 & Veto & Sonnet & 1 & 182 & 888 & 871 & 17 & 0 & 157 & 552 & 0 \\
64 & Veto & Sonnet & 2 & 184 & 924 & 891 & 33 & 0 & 184 & 544 & 0 \\
64 & Veto & Opus & 1 & 182 & 1029 & 1019 & 10 & 0 & 158 & 696 & 0 \\
64 & Veto & Opus & 2 & 184 & 1051 & 1035 & 16 & 0 & 133 & 736 & 0 \\
64 & Veto & GPT-6 Luna & 1 & 182 & 1508 & 1471 & 37 & 0 & 166 & 1144 & 0 \\
64 & Veto & GPT-6 Luna & 2 & 184 & 1480 & 1447 & 33 & 0 & 162 & 1120 & 0 \\
\bottomrule
\end{tabular}

\end{table*}

\end{document}
